% TOP_PROBLEM: {"id": 2637, "title": "Kourovka Notebook Problem 21.128", "category_id": 17, "category": {"id": 17, "name": "group_theory", "display_name": "Group Theory", "description": "Problems about groups, group actions, representations, and related algebraic structures.", "slug": "group-theory", "order_index": 17, "created_at": "2026-07-31T15:26:25.671Z"}, "status": "open", "problem_number": "KOU-21.128", "set_id": 10}
% FIRST_POSTED: 2026-10-01
% ENTRY_KIND: statement_only
% UnsolvedMath ID 2637; source catalogue code KOU-21.128
% !TeX program = lualatex
% Definition and sources only; this file is not an LLM solution attempt.
\documentclass[11pt,a4paper]{article}
\usepackage[margin=25mm]{geometry}
\usepackage{fontspec}
\setmainfont{lmroman10-regular.otf}[BoldFont=lmroman10-bold.otf,
 ItalicFont=lmroman10-italic.otf,BoldItalicFont=lmroman10-bolditalic.otf]
\usepackage{amsmath,amssymb,mathtools}
\usepackage{unicode-math}
\setmathfont{latinmodern-math.otf}
\AtBeginDocument{\let\setminus\smallsetminus}
\usepackage{xurl,hyperref}
\hypersetup{colorlinks=true,urlcolor=blue}
\setcounter{secnumdepth}{0}
\setlength{\parindent}{0pt}
\setlength{\parskip}{5pt}
\setlength{\emergencystretch}{3em}
\begin{document}
\section*{Kourovka Notebook Problem 21.128}\label{2637}
{\small UnsolvedMath ID 2637 · KOU-21.128 · Group Theory · Kourovka Notebook - New Problems, Issue 21}
\subsection{Definitions and mathematical statement}
For two symbols $x,y$ and an integer $m\ge2$, let $[x,y]_m$ denote the alternating positive word of length $m$ beginning with $x$; for example $[x,y]_3=xyx$ and $[x,y]_4=xyxy$. This notation denotes a word, not a group commutator.

For integers $a,b,c\ge2$, define
\[
 A(a,b,c)=\left\langle s_1,s_2,s_3,s_4\ \middle|\
 \begin{array}{l}
 [s_1,s_2]_a=[s_2,s_1]_a,\\
 {}[s_2,s_3]_b=[s_3,s_2]_b,\\
 {}[s_3,s_4]_c=[s_4,s_3]_c,\\
 s_is_j=s_js_i\quad(|i-j|\ge2)
 \end{array}\right\rangle.
\]
The spherical Artin groups in the question are $A[F_4]=A(3,4,3)$ and $A[H_4]=A(5,3,3)$. Their generators are not required to have order two; adding such relations would instead give the associated finite Coxeter groups.

Two groups $G_1,G_2$ are abstractly commensurable if there are finite-index subgroups $H_i\le G_i$ and an isomorphism $H_1\cong H_2$. Are $A[F_4]$ and $A[H_4]$ commensurable in this sense?

The subgroup indices are positive finite integers but need not be equal. The witnessing subgroups need not be normal, and the isomorphism is an abstract group isomorphism, with no requirement to preserve any chosen generating system. Consequently the question is about these two explicitly presented infinite Artin groups, rather than about an equality of their Coxeter diagrams or their finite Coxeter quotients.
\subsection{Short English statement}
Do the spherical Artin groups of types $F_4$ and $H_4$ have isomorphic finite-index subgroups?
\subsection{Sources}
\begin{itemize}
\item[S1] ULAM.AI and the UnsolvedMath Contributors, supplied problems.json, record 2637 (KOU-21.128), Kourovka Notebook - New Problems, Issue 21. Catalogue identity and source transcription; CC BY 4.0.
\url{https://huggingface.co/datasets/ulamai/UnsolvedMath}.
\item[S2] The Kourovka Notebook, 21st edition, version 47 (30 September 2026), new problems, Problem 21.128. Expanded formulation of the supplied catalogue statement.
\url{https://arxiv.org/pdf/1401.0300v47}.
\item[S3] I. Soroko, Artin groups of types F4 and H4 are not commensurable with that of type D4, Topology Appl. 300 (2021), 107770; standard Artin presentations and the separate F4/H4 comparison.
\url{https://arxiv.org/abs/2009.10471}.
\end{itemize}
\end{document}
